\chapter{Methodology}
\label{ch:method}

This chapter describes the dataset, the pre-processing, the network architecture,
the training objective and the evaluation protocol. The complete pipeline is
implemented in PyTorch \citep{paszke2019pytorch} using the MONAI medical-imaging
library \citep{cardoso2022monai}, and every experiment is driven from a single
configurable entry point with a fixed random seed.

\section{Dataset}
\label{sec:data}

We use the BraTS~2020 training dataset \citep{menze2015brats,bakas2017advancing},
which contains 369 subjects. Each subject provides four co-registered MRI
modalities --- native T1, contrast-enhanced T1 (T1Gd), T2 and T2-FLAIR --- together
with an expert segmentation. All volumes are skull-stripped, rigidly registered to
a common anatomical template and resampled to $1\,\mathrm{mm}^3$ isotropic
resolution, giving a spatial size of $240\times240\times155$ voxels per modality.

Each subject is labelled by grade in the dataset's metadata. Of the 369 subjects,
293 are high-grade glioma (HGG, i.e.\ glioblastoma) and 76 are lower-grade glioma
(LGG). Because this project concerns glioblastoma specifically, we filter the
dataset to the \textbf{293 HGG subjects} and use only these for training and
evaluation. This keeps the clinical framing of the study accurate.

The ground-truth annotation assigns each voxel one of four classes. The raw BraTS
labels are $\{0,1,2,4\}$; following common practice we remap the enhancing-tumour
label $4\rightarrow 3$ so that the classes are contiguous:

\begin{table}[h]
\centering
\caption{Segmentation classes used in this work (after label remapping).}
\label{tab:classes}
\begin{tabular}{@{}clc@{}}
\toprule
Index & Region & Raw BraTS label \\
\midrule
0 & Background / healthy tissue & 0 \\
1 & Necrotic and non-enhancing tumour core (NCR/NET) & 1 \\
2 & Peritumoral edema (ED) & 2 \\
3 & Enhancing tumour (ET) & 4 \\
\bottomrule
\end{tabular}
\end{table}

The 293 subjects are split into training and validation sets in an 85:15 ratio
using a fixed random generator (seed 42), giving 249 training and 44 validation
subjects. The split is deterministic so that results are reproducible.

\section{Pre-processing and augmentation}
\label{sec:preproc}

MRI intensities are not calibrated in absolute units and vary between scanners, so
each modality of each subject is independently \emph{z-score normalised} to zero
mean and unit variance over the brain volume:
\begin{equation}
\hat{x}_{c} = \frac{x_{c} - \mu_{c}}{\sigma_{c}},
\end{equation}
where $x_c$ is the intensity of channel (modality) $c$ and $\mu_c,\sigma_c$ are its
mean and standard deviation. The four normalised modalities are stacked into a
single input tensor of shape $4\times D\times H\times W$.

Because a full $240\times240\times155$ volume with a deep 3D network exceeds the
memory of a single consumer GPU, each volume is cropped to a fixed patch of
$128\times128\times128$ voxels centred on the brain. During training only, we
apply light on-the-fly data augmentation to improve generalisation: random flips
along each spatial axis and random $90^{\circ}$ rotations in the axial plane, each
applied with probability $0.5$. A final pad-or-crop step guarantees that every
augmented sample retains the exact patch size so that samples can be batched. No
augmentation is applied at validation time.

\section{Network architecture}
\label{sec:arch}

The segmentation network is a three-dimensional U-Net
\citep{ronneberger2015unet,cicek20163dunet}. It follows the standard
encoder--decoder design: the encoder progressively downsamples the input while
increasing the number of feature channels through the sequence
$(32, 64, 128, 256, 512)$, using strided convolutions with a stride of two at each
of the four downsampling stages; the decoder mirrors this path, upsampling back to
the input resolution. Skip connections concatenate encoder features with the
corresponding decoder features at each resolution, allowing the network to combine
high-level semantic context with fine spatial detail. Each stage uses residual
convolutional units with batch normalisation and a dropout rate of $0.2$ for
regularisation. The network takes the four-channel input volume and outputs four
per-voxel class logits. In this configuration the model has approximately
19.2 million trainable parameters.

\section{Training objective}
\label{sec:loss}

Brain tumour segmentation is highly class-imbalanced: the tumour sub-regions
occupy a small fraction of the volume relative to background. We therefore combine
a region-overlap loss with a per-voxel classification loss, using the sum of the
soft Dice loss and cross-entropy \citep{milletari2016vnet,sudre2017generalised}.
For a class $k$, let $p_{i,k}\in[0,1]$ be the predicted (softmax) probability at
voxel $i$ and $g_{i,k}\in\{0,1\}$ the one-hot ground truth. The soft Dice loss
averaged over the $K$ classes is
\begin{equation}
\mathcal{L}_{\mathrm{Dice}} = 1 - \frac{1}{K}\sum_{k=1}^{K}
\frac{2\sum_{i} p_{i,k}\,g_{i,k} + \epsilon}
     {\sum_{i} p_{i,k} + \sum_{i} g_{i,k} + \epsilon},
\end{equation}
where $\epsilon$ is a small constant that avoids division by zero. The
cross-entropy term is
\begin{equation}
\mathcal{L}_{\mathrm{CE}} = -\frac{1}{N}\sum_{i=1}^{N}\sum_{k=1}^{K}
g_{i,k}\,\log p_{i,k},
\end{equation}
and the total objective is $\mathcal{L} = \mathcal{L}_{\mathrm{Dice}} +
\mathcal{L}_{\mathrm{CE}}$. The Dice term directly optimises overlap and is robust
to imbalance, while the cross-entropy term provides stable, well-behaved gradients
for per-voxel classification.

\section{Optimisation and training protocol}
\label{sec:training}

The network is optimised with AdamW \citep{loshchilov2019adamw} at an initial
learning rate of $10^{-4}$ and weight decay $10^{-4}$. The learning rate follows a
cosine-annealing schedule \citep{loshchilov2017sgdr} over the course of training.
Training uses a batch size of two $128^3$ patches and automatic mixed precision to
reduce memory use and accelerate computation on the GPU. We train for up to 80
epochs with early stopping (patience 15) on the validation mean Dice, keeping the
checkpoint with the highest validation score. All runs use a fixed seed for
reproducibility, and experiments were carried out on a single NVIDIA GeForce
RTX~5050 laptop GPU with 8~GB of memory.

\section{Evaluation metrics}
\label{sec:metrics}

We report two complementary metrics per sub-region. The \textbf{Dice similarity
coefficient} measures volumetric overlap between the predicted set of voxels $P$
and the ground-truth set $G$:
\begin{equation}
\mathrm{Dice}(P,G) = \frac{2\,|P \cap G|}{|P| + |G|} \in [0,1],
\end{equation}
with higher values indicating better agreement. Because Dice is insensitive to
the exact shape of the boundary, we also report the \textbf{95th-percentile
Hausdorff distance} (HD95), which measures boundary error and is more sensitive to
outliers. With $d(a,B)=\min_{b\in B}\lVert a-b\rVert$ the distance from a point to
a set, the 95th-percentile Hausdorff distance is
\begin{equation}
\mathrm{HD95}(P,G) = \max\Big\{
\operatorname*{P_{95}}_{a\in \partial P} d(a,\partial G),\;
\operatorname*{P_{95}}_{b\in \partial G} d(b,\partial P)
\Big\},
\end{equation}
where $\partial P,\partial G$ are the surfaces of the prediction and ground truth
and $\mathrm{P}_{95}$ denotes the 95th percentile. HD95 is reported in millimetres
(equivalently voxels, since the data are $1\,\mathrm{mm}^3$ isotropic), and lower
values are better. We report both metrics for each foreground sub-region and their
mean, and we visualise predictions as overlays on the FLAIR modality next to the
ground truth.
